\section{Results}

\subsection{Task Compression Clusters Exist (RQ1)}

Gap statistic analysis reveals three distinct compression response clusters. The gap criterion is satisfied: $\text{Gap}(3) = 0.957 \geq \text{Gap}(4) - s_4 = 0.898$.

\begin{table}[h]
\centering
\caption{Gap Statistic Results}
\begin{tabular}{@{}lll@{}}
\toprule
k & Gap Value & Standard Error \\
\midrule
1 & -0.078 & 0.099 \\
2 & 0.713 & 0.101 \\
3 & \textbf{0.957} & \textbf{0.119} \\
4 & 1.018 & 0.120 \\
\bottomrule
\end{tabular}
\end{table}

The three clusters align with interpretable task categories:
\begin{itemize}
    \item \textbf{Cluster 0 (Multi-doc + Code):} High eviction sensitivity
    \item \textbf{Cluster 1 (Single-doc + Few-shot):} Moderate sensitivity
    \item \textbf{Cluster 2 (Summarization + Synthetic):} Different compression profile
\end{itemize}

Silhouette score is 0.411, indicating reasonable but imperfect cluster structure.

\subsection{Attention Entropy Discriminates Task Domains (RQ2)}

One-way ANOVA reveals highly significant entropy differences across domains:

\begin{table}[h]
\centering
\begin{tabular}{@{}ll@{}}
\toprule
Metric & Value \\
\midrule
F-statistic & 38.05 \\
p-value & $2.92 \times 10^{-26}$ \\
Effect size ($\eta^2$) & 0.522 \\
\bottomrule
\end{tabular}
\end{table}

The effect size $\eta^2 = 0.522$ means 52\% of entropy variance is explained by task domain---a large effect.

\subsection{Entropy-Tolerance Relationship (RQ3 --- Inconclusive)}

The hypothesis that high-entropy tasks tolerate eviction better could not be validated (p=0.326) due to measurement limitations: baseline accuracy of 6.7\% made retention measurement unreliable. This identifies an evaluation gap for future work.
